\scanpage{397}\mspage{136}

Par contre, (*) est iso quand $C$ est une catégorie de Thomason-Kan. Il suffit de voir que pour \deleted{tout} $Z$ de $(\mathrm{Cat})$, le $\mathrm{Hom}$ de $\mathrm{hot}(Z)$ dans les deux membres est \unsure{isomorphé} par la flèche déduite de (*). Or si $C$ est T.K., de même $\underline{\mathrm{Hom}}(Y,C)$ (immédiat sur la définition \inserted{de T.K.} en termes de prolongements de \deleted{\unsure{Y}}

\diagram{\begin{tikzcd}[ampersand replacement=\&,row sep=small,column sep=small]
A_0\arrow[dr]\&\& A\arrow[ll,"i"']\arrow[dl,dashed]\&\ldots\\
\& C\&\&
\end{tikzcd}}

), donc

% slanted
\margin{un peu vif ! Il faut savoir que les applications \unsure{cofibrantes|cofibrations} \inserted{\unsure{$i$|?}} de Thomason le restent après multiplication cartésienne par $Y$ \ldots}

\[
\begin{aligned}
&\mathrm{Hom}_{\mathrm{Hot}}(\mathrm{hot}(Z),\mathrm{hot}(\underline{\mathrm{Hom}}(Y,C)))\\
&\xleftarrow{\sim}\pi_0\big(\underbrace{\underline{\mathrm{Hom}}_{\mathrm{Cat}}(Z,\underline{\mathrm{Hom}}(Y,C))}_{\underline{\mathrm{Hom}}_{\mathrm{Cat}}(Z\times Y,C)}\big)
\simeq\big(\pi_0\underline{\mathrm{Hom}}(Z\times Y,C)\big),
\end{aligned}
\]

et \deleted{$\mathrm{Hom}(\mathrm{hot}(Z),$}

\[
\begin{aligned}
&\mathrm{Hom}_{\mathrm{Hot}}\big(\mathrm{hot}(Z),\underline{\mathrm{Hom}}_{\mathrm{Hot}}(\mathrm{hot}(\unsure{Y}),\mathrm{hot}(C))\big)
\simeq\mathrm{Hom}_{\mathrm{Hot}}\big(\underbrace{\mathrm{hot}(Z)\times\mathrm{hot}(Y)}_{\mathrm{hot}(Z\times Y)},\mathrm{hot}(C)\big)\\
&\underset{\text{\inserted{car $C$ est T.K.}}}{\xleftarrow{\sim}}\pi_0\big(\underline{\mathrm{Hom}}_{\mathrm{Cat}}(Z\times Y,C)\big),
\end{aligned}
\]

Ok.

On voit donc que les catégories $\mathrm{Ch}_{\infty,\overline\Sigma}(X;a,b)$ n'ont aucune tendance à être des catégories de Thomason-Kan. Pourtant, \unsure{son|ses} type dans $\mathrm{Hot}$, doit être ``le bon'', sans avoir d'abord à remplacer $X$ par une enveloppe de Th.K. $\overline X$.

% archive sheet number: 396
